% Propietario conservado: /Users/ruben/Documents/New project/output/EXTREMA_MEDIA_RAZON_RECIPROCIDAD_ES_EN_20260918/ARTICULO/ES/source/nuclear/08.tex
\chapter{Conservation des formes d’action et de l’aire orientée dans les régimes tritiques}
\label{x_reciprocidad:nuclear:8}
\section{Bilan entre deux formes d’action}

Les sorties angulaires \(A,C^*\), dans une même unité, déterminent la forme \(\eta=\operatorname{artanh}(C^*/A)\) sur le domaine \(A\ne0\), \(|C^*/A|<1\). Deux formes successives \(\eta_n,\eta_{n+1}\) étant reçues, nous définissons

\[
M_n=\operatorname{diag}(e^{\eta_n/2},e^{-\eta_n/2}),\qquad
G_n=M_n^{-\mathsf T}M_n^{-1},\qquad
I_n(z)=\tfrac12z^\mathsf TG_nz.
\]

Ici, \(z=(Q,P)^\mathsf T\), avec les deux coordonnées de dimension racine d’action. L’échelle \(\hbar\), lorsqu’elle fixe le contour \(I_n=\hbar\), est la sortie d’action précédemment construite ; elle n’intervient pas comme cible d’ajustement.

Soit \(R_n\) la rotation canonique de l’avancée angulaire \(\theta_n\). Le transport sans avancée angulaire et le transport avec cette avancée sont, respectivement,

\begin{equation}
B_n=M_{n+1}M_n^{-1},\qquad
C_n=M_{n+1}R_nM_n^{-1}.
\label{x_reciprocidad:nuc08:eq:1}
\end{equation}

Tous deux vont de la carte \(n\) à la carte \(n+1\). Nous définissons la lecture comprimée et son registre complémentaire normalisé :

\begin{equation}
T_n=\frac{8B_n+C_n}{9},\qquad
D_n=\frac{\sqrt8}{9}(C_n-B_n).
\label{x_reciprocidad:nuc08:eq:2}
\end{equation}

L’opérateur \(D_n\) est une réalisation isométriquement normalisée de l’unique motif de différences entre les neuf phases. Le registre de neuf composantes est \((8z,-z,\ldots,-z)/27\), avec \(z=(C_n-B_n)v\) ; la somme de leurs produits scalaires donne le même facteur \(8/81\). Le registre de neuf phases est maintenu lorsque les étiquettes individuelles interviennent dans d’autres lecteurs.

\begin{theorem}
L’action est conservée entre lecture et mémoire :

\begin{equation}
\boxed{G_n=T_n^\mathsf TG_{n+1}T_n+D_n^\mathsf TG_{n+1}D_n.}
\label{x_reciprocidad:nuc08:eq:3}
\end{equation}
\end{theorem}

\begin{proof}
La normalisation des cartes donne
\(B_n^\mathsf TG_{n+1}B_n=G_n=C_n^\mathsf TG_{n+1}C_n\).
En développant \eqref{x_reciprocidad:nuc08:eq:2}, les contributions croisées ont pour coefficients \(+8\) et \(-8\). Les coefficients diagonaux restants sont \(72\) pour \(B_n\) et \(9\) pour \(C_n\), divisés par \(81\). La somme est \(G_n\).
\end{proof}

Avec \(z_{n+1}=T_nz_n\), \(m_n=D_nz_n\), la forme scalaire est

\begin{equation}
\boxed{I_n(z_n)=I_{n+1}(z_{n+1})+I_{n+1}(m_n).}
\label{x_reciprocidad:nuc08:eq:4}
\end{equation}

La somme télescopique, en retenant chaque mémoire dans la carte où elle a été émise, donne

\begin{equation}
I_0(z_0)=I_N(z_N)+\sum_{n=0}^{N-1}I_{n+1}(m_n).
\label{x_reciprocidad:nuc08:eq:5}
\end{equation}

Le transport \(B_n\) est matériel dans cette composition : le remplacer par l’identité lorsque les formes diffèrent modifie les domaines et peut rompre \eqref{x_reciprocidad:nuc08:eq:3}.

\section{Loi exacte d’accumulation à phase variable}

Dans les coordonnées normalisées \(y_n=M_n^{-1}z_n\), la lecture est
\(y_{n+1}=(8I+R_n)y_n/9\). Pour une rotation plane,

\begin{equation}
\left(\frac{8I+R_n}{9}\right)^\mathsf T
\left(\frac{8I+R_n}{9}\right)=r_n I,\qquad
r_n=1-\frac{32}{81}\sin^2(\theta_n/2).
\label{x_reciprocidad:nuc08:eq:6}
\end{equation}

Par conséquent,

\begin{equation}
I_N(z_N)=I_0(z_0)\prod_{n=0}^{N-1}r_n.
\label{x_reciprocidad:nuc08:eq:7}
\end{equation}

\begin{theorem}
Pour \(I_0(z_0)>0\), le terminal tend vers zéro exactement lorsque

\begin{equation}
\sum_{n\ge0}\sin^2(\theta_n/2)=+\infty.
\label{x_reciprocidad:nuc08:eq:8}
\end{equation}
\end{theorem}

\begin{proof}
Soient \(c=32/81\) et \(u_n=\sin^2(\theta_n/2)\in[0,1]\). On a
\(cu_n\leq-\log(1-cu_n)\leq cu_n/(1-c)\), par intégration de \(1/(1-s)\) entre zéro et \(cu_n\). La somme de ces logarithmes diverge exactement lorsque \(\sum u_n\) diverge. En appliquant le logarithme au produit de \eqref{x_reciprocidad:nuc08:eq:7}, on obtient l’équivalence.
\end{proof}

On conserve \eqref{x_reciprocidad:nuc08:eq:5} dans les deux cas. Une suite d’avancées angulaires peut laisser un terminal positif ou transférer toute l’action au registre, selon la somme \eqref{x_reciprocidad:nuc08:eq:8}. Le changement de forme seul conserve l’action ; la dispersion angulaire détermine sa répartition dans cette réalisation. Les phases \(\theta_n\) doivent provenir du lecteur concret de la trajectoire : la preuve ne sélectionne pas rétrospectivement une suite pour imposer un résultat.

\section{Euler étendu et aire orientée commune}

Les générateurs des trois régimes quadratiques sont

\[
J_{+1}=(2C_{A_2}+I)/\sqrt3,\qquad
J_0=N,\qquad
J_{-1}=(2F_{\rm av}^2-3I)/\sqrt5.
\]

Dans leurs cartes bidimensionnelles, ils sont de trace nulle et satisfont
\(J_\tau^2=-\tau I\). Le Coxeter provient de l’antécédent APP, le nilpotent du transport de seuil et \(F_{\rm av}\) de la descente des modes TPK. Les fonctions exponentielles constituent la réalisation analytique ultérieure :


\[
E_\tau(t)=\exp(tJ_\tau)=c_\tau(t)I+s_\tau(t)J_\tau,
\]

où \((c_\tau,s_\tau)\) sont \((\cos,\sin)\), \((1,t)\) et \((\cosh,\sinh)\), respectivement. La trace nulle donne le déterminant un. Avec

\[
\Omega=\begin{pmatrix}0&1\\-1&0\end{pmatrix},
\]

on obtient \(E_\tau^\mathsf T\Omega E_\tau=\Omega\). Dans des cartes variables, on prend
\(C_n=M_{n+1}E_{\tau_n}(t_n)M_n^{-1}\) et le même \(B_n\) de \eqref{x_reciprocidad:nuc08:eq:1}.

\begin{theorem}
Les trois régimes satisfont

\begin{equation}
\boxed{\Omega=T_n^\mathsf T\Omega T_n+D_n^\mathsf T\Omega D_n,}
\label{x_reciprocidad:nuc08:eq:9}
\end{equation}

avec \(T_n,D_n\) de \eqref{x_reciprocidad:nuc08:eq:2}.
\end{theorem}

\begin{proof}
En dimension deux, \(L^\mathsf T\Omega L=(\det L)\Omega\). Les déterminants de \(M_n\), \(B_n\) et \(C_n\) valent un. Le développement utilisé dans \eqref{x_reciprocidad:nuc08:eq:3}, appliqué à la forme alternée, démontre \eqref{x_reciprocidad:nuc08:eq:9}.
\end{proof}

Dans une carte fixe, les deux déterminants sont

\begin{equation}
\det T_\tau=\frac{65+16c_\tau(t)}{81},\qquad
\det D_\tau=\frac{16(1-c_\tau(t))}{81}.
\label{x_reciprocidad:nuc08:eq:10}
\end{equation}

Cela résulte de \(\det(aI+C)=a^2+a\operatorname{tr}C+\det C\) et \(\operatorname{tr}E_\tau=2c_\tau\). Leur somme est un. Dans les représentants \(C_{A_2}\), \(I+N\), \(F_{\rm av}^2\), les paires sont \((19/27,8/27)\), \((1,0)\), \((89/81,-8/81)\).

Le régime parabolique peut produire une mémoire non nulle de rang un et d’aire nulle. Le régime hyperbolique conserve l’aire orientée au moyen d’un complément de signe opposé. Le régime elliptique permet en outre le bilan positif d’action entre les formes des deux fibres. La forme alternée commune et les formes métriques de chaque régime maintiennent leurs types respectifs.

\subsection{Action intégrée et orientation}

Soit \(\gamma\) une courbe fermée \(C^1\) par morceaux dans le plan canonique. Pour toute matrice réelle \(L\) d’ordre deux,

\[
\oint_{L\gamma}P\,dQ=(\det L)\oint_\gamma P\,dQ.
\]

En effet, en écrivant \(Q'=aQ+bP\), \(P'=cQ+dP\), les intégrales de \(Q\,dQ\) et \(P\,dP\) sont nulles, et \(\oint Q\,dP=-\oint P\,dQ\). Il reste le coefficient \(ad-bc\). En appliquant \eqref{x_reciprocidad:nuc08:eq:10},

\begin{equation}
\boxed{\oint_\gamma P\,dQ
=\oint_{T_n\gamma}P\,dQ+\oint_{D_n\gamma}P\,dQ.}
\label{x_reciprocidad:nuc08:eq:11}
\end{equation}

Les signes sont ceux des courbes transportées. L’identité est conservée pour les trois régimes et pour les changements de forme de déterminant un.

\section{Le facteur commun du lecteur exponentiel}

Le lecteur radional de \cref{x_reciprocidad:iv:ellipse:section,x_reciprocidad:euler-trit-evaluaciones} contient aussi \(e^{-x}\exp(-yJ_\tau)\). Ce facteur est conservé lorsqu’il agit dans le lecteur. Soit

\[
C_n=e^{-x_n}M_{n+1}E_{\tau_n}(t_n)M_n^{-1}.
\]

Maintenant, \(C_n^\mathsf T\Omega C_n=e^{-2x_n}\Omega\). Le même développement donne

\begin{equation}
T_n^\mathsf T\Omega T_n+D_n^\mathsf T\Omega D_n
=\frac{8+e^{-2x_n}}9\Omega.
\label{x_reciprocidad:nuc08:eq:12}
\end{equation}

Pour \(x_n\geq0\), nous définissons un registre complémentaire d’échelle

\[
D_{{\rm esc},n}=\frac{\sqrt{1-e^{-2x_n}}}{3}B_n.
\]

Alors l’identité complète est

\begin{equation}
\Omega=T_n^\mathsf T\Omega T_n+D_n^\mathsf T\Omega D_n
+D_{{\rm esc},n}^\mathsf T\Omega D_{{\rm esc},n}.
\label{x_reciprocidad:nuc08:eq:13}
\end{equation}

C’est une composition définie ici pour représenter explicitement le défaut d’échelle du lecteur. Son coefficient provient de \eqref{x_reciprocidad:nuc08:eq:12}. On n’attribue pas à ce registre, par son nom, une identité avec un autre secteur physique du corpus ; on conserve l’application afin que cette comparaison puisse être faite avec leurs opérateurs et domaines respectifs.

\section{Raffinement d’action sur les émissions complètes}

Pour chaque préfixe enrichi \(h\), soit \(M_h\) sa carte et \(G_h=M_h^{-\mathsf T}M_h^{-1}\). Pour chaque émission survivante \(\epsilon\), soit \(R_\epsilon\) une rotation réalisée et


\[
U_\epsilon=M_{h\epsilon}R_\epsilon M_h^{-1},\qquad
(\mathscr U z)_{h\epsilon}=\sqrt{p_h(\epsilon)}U_\epsilon z_h.
\]

Les poids proviennent de la mesure des émissions de la construction commune ; \(\sum_\epsilon p_h(\epsilon)=1\), avec les multiplicités préservées. De \(U_\epsilon^\mathsf TG_{h\epsilon}U_\epsilon=G_h\), il découle

\begin{equation}
\sum_{h,\epsilon}I_{h\epsilon}((\mathscr U z)_{h\epsilon})
=\sum_h I_h(z_h).
\label{x_reciprocidad:nuc08:eq:14}
\end{equation}

La preuve somme d’abord les poids des descendants d’un parent. Dans une trajectoire, les facteurs \(M_h^{-1}M_h\) intermédiaires s’annulent, et les rotations se composent dans l’ordre enregistré. Cela conserve les étiquettes et la mémoire sans mélanger les enfants de parents différents. La forme alternée en somme directe satisfait la même identité de transport pour les trois régimes symplectiques, en remplaçant \(R_\epsilon\) par \(E_{\tau_\epsilon}(t_\epsilon)\).


Les complétions métriques sont prises avec \(\sum_h z_h^\mathsf TG_hz_h\), de sorte que ces applications sont isométriques. Identifier en outre cette norme à une somme directe euclidienne fixe exige des bornes uniformes de \(M_h\) et \(M_h^{-1}\) ; cette identification supplémentaire n’est pas utilisée dans \eqref{x_reciprocidad:nuc08:eq:14}.
